\chapter{Buffers, Windows, Tabs, Sessions}
\label{ch:buffers}

Chapter~\ref{ch:ontology} named buffers, windows, and tabs as distinct entities and warned against conflating them: a buffer is loaded content, a window is a viewport onto a buffer, and a tab page is a saved arrangement of windows. This chapter returns to that distinction with the actual commands for managing each, plus sessions, which extend the same idea of "an arrangement worth returning to" across the boundary of a single editing session entirely.

\section{Managing Buffers}

\texttt{:ls} (or its synonym \texttt{:buffers}) lists every buffer currently loaded, whether or not it is visible in any window, each entry numbered and flagged with its status: \texttt{\%} for the buffer in the current window, \texttt{\#} for the alternate buffer from Chapter~\ref{ch:marks}, \texttt{a} for active (loaded and visible), \texttt{h} for hidden (loaded but not currently displayed anywhere), and \texttt{+} for a buffer with unsaved modifications. \texttt{:b}\textit{n} switches the current window to buffer number \textit{n}, and \texttt{:b}\textit{name} switches to a buffer by a (partially matched) filename, which is generally the faster of the two in practice since it does not require first consulting \texttt{:ls} to look up a number. \texttt{:bnext} and \texttt{:bprevious} (abbreviated \texttt{:bn} and \texttt{:bp}) cycle through the buffer list in order.

A buffer becomes hidden, rather than closed, when its window is closed while the buffer itself still has unsaved changes Vim is unwilling to discard silently, or when \texttt{:hide} or an equivalent command is used deliberately. A hidden buffer retains its full content and undo history; it is simply not currently displayed in any window. \texttt{:bd} (buffer delete) removes a buffer from the buffer list entirely, which is a distinct and more final operation than merely closing the window that was displaying it, and will refuse to proceed if the buffer has unsaved changes unless forced with \texttt{:bd!}.

\section{Splitting Windows}

\texttt{:split} (or \texttt{:sp}) divides the current window horizontally, creating a new window above the existing one, both displaying the same buffer initially. \texttt{:vsplit} (or \texttt{:vsp}) divides vertically instead. Either command can be given a filename, \texttt{:split file.txt}, to open a different file in the new window directly rather than duplicating the current buffer into it. \texttt{Ctrl-w s} and \texttt{Ctrl-w v} are the Normal-mode equivalents, useful for splitting without dropping into Command-line mode at all.

Once split, \texttt{Ctrl-w} followed by a direction key (\texttt{h}, \texttt{j}, \texttt{k}, or \texttt{l}) moves the cursor to the window in that direction, the same four keys used for character-wise motion in Chapter~\ref{ch:coordinates}, repurposed here for window-wise movement, consistent with the ergonomic argument from Chapter~\ref{ch:different} that the home row should serve as many purposes as reasonably possible rather than requiring the hand to travel for a conceptually related but differently scaled operation. \texttt{Ctrl-w w} cycles to the next window regardless of direction, and \texttt{Ctrl-w c} closes the current window (distinct from \texttt{:bd}, since closing a window does not necessarily delete the buffer it was showing, per the hidden-buffer behavior above).

\section{Resizing and Rearranging}

\texttt{Ctrl-w =} resizes all windows to be as equal in size as the current layout permits, a useful reset after manual resizing has left the screen unevenly divided. \texttt{Ctrl-w +} and \texttt{Ctrl-w -} grow and shrink the current window's height by one line (prefixed with a count for a larger adjustment, following the same count convention established in Chapter~\ref{ch:language}); \texttt{Ctrl-w <} and \texttt{Ctrl-w >} do the same for width. \texttt{Ctrl-w r} rotates the windows in the current layout, and \texttt{Ctrl-w x} exchanges the current window with the next one, both useful for correcting a split that was created in a less convenient order than intended without having to close and reopen anything.

\section{Tabs}

A tab page, as established in Chapter~\ref{ch:ontology}, holds an entire arrangement of windows, not a single file. \texttt{:tabnew} (optionally given a filename) opens a new tab; \texttt{gt} and \texttt{gT} move to the next and previous tab respectively, and \texttt{n gt} moves directly to tab number \textit{n}. \texttt{:tabclose} closes the current tab, and \texttt{:tabonly} closes every tab except the current one. A count prefixed to \texttt{:tabnew}, or the command \texttt{:tabnew} issued repeatedly, has no relationship to the count-prefixed motions of Chapter~\ref{ch:coordinates}; the parallel keystroke vocabulary (\texttt{g} followed by a letter, echoing \texttt{g;} and \texttt{g,} from the change list in Chapter~\ref{ch:marks}) is a naming convention rather than a grammatical extension of the operator-motion rule from Part II.

The practical case for tabs over simply splitting one very large window many times is organizational: a tab groups a set of windows around a single coherent task (one project's relevant files, arranged in whatever split configuration suits that task), and switching tabs switches that entire grouping at once, rather than requiring individual windows to be closed and reopened as the task changes. A reader working on two unrelated pieces of a project simultaneously will generally find two tabs, each internally split as needed, more legible than one tab split many ways across both tasks at once.

\section{Sessions}

Everything covered in this chapter so far, the buffer list, the window layout, the tab arrangement, and each window's cursor position and view, can be captured as a unit and restored later. \texttt{:mksession} (optionally given a filename, defaulting to \texttt{Session.vim}) writes the current arrangement to a Vimscript file; running Vim with \texttt{vim -S} \textit{filename}, or sourcing that file from within a running session with \texttt{:source}, restores the arrangement exactly, reopening every buffer, recreating every split and tab, and repositioning every cursor.

Sessions are the natural endpoint of the buffer, window, and tab material in this chapter, because they answer the same "how do I get back to this without reconstructing it by hand" question that motivated marks and registers in the two preceding chapters, applied now to an entire working environment rather than a single position or fragment of text. A reader working across several distinct projects, each with its own characteristic arrangement of files and splits, can maintain a separate session file per project and restore the full working context for any of them with a single command, rather than manually reopening and re-splitting the same handful of files each time work resumes.

\section{Buffers, Windows, Tabs, and Sessions as a Single Hierarchy}

It is worth closing Part IV by restating the hierarchy this chapter has covered, since it is easy to lose track of amid the individual commands: a session contains tabs, a tab contains windows, and a window displays a buffer, which may or may not correspond to a file on disk, per the distinction drawn in Chapter~\ref{ch:ontology}. Each level of this hierarchy has its own management commands, largely independent of the levels above and below it, and a reader disoriented by an unfamiliar Vim session (an unexpected number of open windows, an unfamiliar tab arrangement) can generally resolve that disorientation by asking, in order, which buffer is currently active, which window it is displayed in, and which tab that window belongs to, rather than treating the whole arrangement as an undifferentiated puzzle to be solved at once.
